% from ega2-1-fr.tex
\oldpage[II]{22}
linéaires, à coefficients dans $S_0$, des éléments
$f_1^{\alpha_1}\cdots f_r^{\alpha_r}x_j$ tels que les $\alpha_i$ soient des
entiers $\geq 0$ vérifiant $k_j+\sum_i\alpha_i h_i=n$~; pour chaque $j$, il
n'y a qu'un nombre fini de systèmes $(\alpha_i)$ vérifiant cette équation,
puisque les $h_i$ sont $>0$, d'où la première assertion de (i)~; la seconde
est évidente. Soit d'autre part $h$ le p.p.c.m. des $h_i$ et posons
$g_i=f_i^{h/h_i}$ ($1\leq i\leq r$) de sorte que tous les $g_i$ sont de
degré $h$~; soient $z_\mu$ les éléments de $M$ de la forme
$f_1^{\alpha_1}\cdots f_r^{\alpha_r}x_j$ avec
$0\leq\alpha_i<h/h_i$ pour $1\leq i\leq r$~; ces éléments sont en nombre
fini, soit $n_1$ le plus grand de leurs degrés. Il est clair que pour
$n\geq n_1$, tout élément de $M_{n+h}$ est combinaison linéaire des $z_\mu$
dont les coefficients sont des monômes de degré $>0$ par rapport aux $g_i$,
donc on a $M_{n+h}=S_hM_n$, ce qui établit (ii). De la même manière, on voit
(pour tout $d>0$) qu'un élément de $M^{(d,k)}$ est combinaison linéaire, à
coefficients dans $S_0$, d'éléments de la forme
$g^df_1^{\alpha_1}\cdots f_r^{\alpha_r}x_j$ avec
$0\leq\alpha_i<d$, $g$ étant un élément homogène de $S$~; d'où (iii)~; (iv)
résulte alors de (iii) et du lemme (\hyperref[II.2.1.3-fr]{2.1.3}), en
prenant $M=S_+$, puisque $(S_+)^{(d)}=(S^{(d)})_+$. L'assertion de (v) se
déduit de (ii) en prenant $M=S$. Enfin, pour un $n$ donné, il n'y a qu'un
nombre fini de systèmes $(\alpha_i)$ tels que $\alpha_i\geq 0$ et
$\sum_i\alpha_i<n$, donc si $m_0$ est la plus grande valeur de la somme
$\sum_i\alpha_i h_i$ pour ces systèmes, on a $S_m\subset S_+^n$ pour
$m>m_0$, ce qui prouve (vi).
\end{proof}

\begin{corollary}[2.1.7]
\label{II.2.1.7-fr}
Si $S$ est noethérien, il en est de même de $S^{(d)}$ pour tout entier
$d>0$.
\end{corollary}

\begin{proof}
Cela résulte de (\hyperref[II.2.1.5-fr]{2.1.5}) et de
(\hyperref[II.2.1.6-fr]{2.1.6, (iv)}).
\end{proof}

\begin{env}[2.1.8]
\label{II.2.1.8-fr}
Soit $\mathfrak{p}$ un idéal premier gradué de l'anneau gradué $S$~;
$\mathfrak{p}$ est donc somme directe des sous-groupes
$\mathfrak{p}_n=\mathfrak{p}\cap S_n$. Supposons que $\mathfrak{p}$ ne
contienne pas $S_+$. Alors, si $f\in S_+$ n'appartient pas à
$\mathfrak{p}$, la relation $f^nx\in\mathfrak{p}$ est équivalente à
$x\in\mathfrak{p}$~; en particulier, si $f\in S_d$ ($d>0$), pour tout
$x\in S_{m-nd}$, la relation $f^nx\in\mathfrak{p}_m$ équivaut à
$x\in\mathfrak{p}_{m-nd}$.
\end{env}

\begin{proposition}[2.1.9]
\label{II.2.1.9-fr}
Soit $n_0$ un entier $>0$~; pour tout $n\geq n_0$ soit $\mathfrak{p}_n$ un
sous-groupe de $S_n$. Pour qu'il existe un idéal premier gradué
$\mathfrak{p}$ de $S$ ne contenant pas $S_+$ et tel que
$\mathfrak{p}\cap S_n=\mathfrak{p}_n$ pour tout $n\geq n_0$, il faut et il
suffit que les conditions suivantes soient vérifiées~:
\begin{enumerate}
  \item[$1^\circ$] $S_m\mathfrak{p}_n\subset\mathfrak{p}_{m+n}$ pour tout
    $m\geq 0$ et tout $n\geq n_0$.
  \item[$2^\circ$] Pour $m\geq n_0$, $n\geq n_0$, $f\in S_m$, $g\in S_n$,
    la relation $fg\in\mathfrak{p}_{m+n}$ entraîne
    $f\in\mathfrak{p}_m$ ou $g\in\mathfrak{p}_n$.
  \item[$3^\circ$] $\mathfrak{p}_n\neq S_n$ pour un $n\geq n_0$ au moins.
\end{enumerate}
En outre l'idéal premier gradué $\mathfrak{p}$ est alors unique.
\end{proposition}

\begin{proof}
Il est évident que les conditions $1^\circ$ et $2^\circ$ sont nécessaires.
En outre, si $\mathfrak{p}\not\supset S_+$, il existe au moins un $k>0$ tel que
$\mathfrak{p}\cap S_k\neq S_k$~; si $f\in S_k$ n'appartient pas à
$\mathfrak{p}$, la relation $\mathfrak{p}\cap S_n=S_n$ entraîne
$\mathfrak{p}\cap S_{n-mk}=S_{n-mk}$ d'après
(\hyperref[II.2.1.8-fr]{2.1.8})~; donc, si
$\mathfrak{p}\cap S_n=S_n$ à partir d'une certaine valeur de $n$, on aurait
$\mathfrak{p}\supset S_+$ contrairement à l'hypothèse, ce qui prouve que
$3^\circ$ est nécessaire. Inversement, supposons satisfaites les conditions
$1^\circ$, $2^\circ$ et $3^\circ$. Notons que, si pour un entier
$d\geq n_0$, $f\in S_d$ n'appartient pas à $\mathfrak{p}_d$, alors, si
$\mathfrak{p}$ existe, $\mathfrak{p}_m$, pour $m<n_0$, est nécessairement
égal à l'ensemble des $x\in S_m$ tels que
$f^rx\in\mathfrak{p}_{m+rd}$, sauf pour un nombre fini de valeurs de $r$.
Cela prouve déjà que si $\mathfrak{p}$ existe, il est unique. Reste à montrer
que si on définit les $\mathfrak{p}_m$ pour $m<n_0$ par la condition
précédente, $\mathfrak{p}=\sum_{n=0}^{\infty}\mathfrak{p}_n$ est un idéal
premier. Notons d'abord qu'en vertu de $2^\circ$, pour $m\geq n_0$,
$\mathfrak{p}_m$ est aussi défini comme l'ensemble des $x\in S_m$ tels que
$f^rx\in\mathfrak{p}_{m+rd}$, sauf pour un nombre fini de valeurs de $r$.
Cela
